\section{Linux Networking}\label{linux-networking}

Per approfondire meglio il discorso relativo al routing dei pacchetti e
delle interfacce di rete, si vuole elencare ad alto livello cosa succede
all'interno del kernel linux e quali sono gli oggetti fondamentali.

\subsection{Network Device Driver}\label{network-device-driver}

Un network device opera sempre a livello 2 o 3 dello stack di rete, ed è
rappresentato nel kernel dalla struttura
\passthrough{\lstinline!net\_device!}.\footnote{\label{fn:ref_23}R. Rosen, \emph{Linux
  Kernel Networking: Implementation and Theory}, Apress, 2014.} Questa
struttura dati contiene varie informazioni:

\begin{itemize}
\tightlist
\item
  \passthrough{\lstinline!IRQ!}: numero identificativo del dispositivo
\item
  \passthrough{\lstinline!MTU!}
\item
  \passthrough{\lstinline!MAC!}: MAC address
\item
  \passthrough{\lstinline!Nome!}: \passthrough{\lstinline!eth0!},
  \passthrough{\lstinline!enp0s3!}, \passthrough{\lstinline!wlan0!},
  ecc\ldots{}
\item
  \passthrough{\lstinline!net\_device\_ops!}: contiene i puntatori alle
  funzioni per la gestione del dispositivo (come le primitive start,
  stop, transmit, receive)
\end{itemize}

Un device driver si occupa di:

\begin{itemize}
\tightlist
\item
  \textbf{Ricevere} i pacchetti destinati al local host, inoltrandoli al
  livello di rete, e successivamente al livello di trasporto.
\item
  \textbf{Trasmettere} i pacchetti generati dal local host, o fa il
  forward dei pacchetti ricevuti verso un'altra interfaccia.
\end{itemize}

Ogni pacchetto in transito in tutti e due i sensi scatena un lookup nel
\emph{sottosistema di routing}, che decide cosa bisogna fare con quel
pacchetto. La fase di routing si distingue in 5 fasi, detti
\textbf{hook}, dove possiamo registrare dei \textbf{callback} per il
\textbf{sottosistema netfilter}.

Per esempio l'hook \passthrough{\lstinline!NF\_INET\_PRE\_ROUTING!},
registra la chiamata al callback per essere eseguito prima che il lookup
del sottosistema di routing venga effettuato. Il callback ritorna un
verdetto sul filtraggio del pacchetto, decidendo se bisogna accettarlo
(\passthrough{\lstinline!NF\_ACCEPT!}), oppure scartarlo
(\passthrough{\lstinline!NF\_DROP!}).

\begin{quote}
Con \passthrough{\lstinline!nftables!} applichiamo le regole di
filtraggio dei pacchetti in modo da intercettare e manipolare il
sottosistema di routing
\end{quote}

\begin{quote}
Oltre ai callback registrati, altri fattori influenzano il modo in cui
il routing viene effettuato. Per esempio un pacchetto viene scartato
quando il TTL viene diminuito e arriva a 0.
\end{quote}

La struttura \passthrough{\lstinline!sk\_buff!} (Socket Buffer)
definisce come un pacchetto nel kernel Linux viene rappresentato, tale
struttura viene manipolata unicamente attraverso la SKB API.

La Socket Buffer viene allocata dinamicamente quando il pacchetto viene
ricevuto sul cavo, ed è la struttura dati utilizzata per manipolare il
pacchetto nel kernel sui vari layer di rete.

All'interno del kernel la struttura dati responsabile per il
sottosistema di routing è la \passthrough{\lstinline!fib\_table!}, ossia
la tabella che raccoglie le rotte per il traffico destinato ad una certa
sottorete.

\subsection{Bridge}\label{bridge}

Un bridge è un dispositivo L2 che capisce il protocollo del livello di
collegamento ed effettua l'inoltro dei dati copiandoli frame per
frame.\footnote{\label{fn:ref_24}C. Benvenuti, \emph{Understanding Linux Network
  Internals}, O'Reilly Media, 2006.} Per farlo deve mantenere nel suo
buffer almeno un frame per volta.

Un bridge assegna un indirizzo al livello di collegamento ad ognuna
delle sue interfacce ed effettua il forward verso di queste se il bridge
non è il destinatario del pacchetto. Opzionalmente può implementare il
livello di rete (Layer 3) per offrire visibilità sulla WAN alle
interfacce collegate. Un bridge viene di solito utilizzato per unire due
LAN differenti, o per creare una LAN.

Un bridge nel kernel Linux è un device virtuale, finché non viene
collegato ad altri device non può né ricevere né trasmettere frame. Usa
il concetto di enslaving: ossia gli vengono assegnate delle interfacce
dette \textbf{enslaved} da controllare.

Quando si tenta di trasmettere su un bridge (es:
\passthrough{\lstinline!br0!}), il device driver gestisce i device
\textbf{enslaved} effettuando il forwarding dei frame.

\subsection{Netfilter}\label{netfilter}

È un sottosistema del kernel che fornisce le seguenti funzionalità
all'interno del sottosistema di routing:\footnote{\label{fn:ref_25}The Netfilter Project,
  \emph{Netfilter Architecture and Packet Traversal Documentation},
  https://www.netfilter.org/} \footnote{\label{fn:ref_26}Netfilter Core Team,
  \emph{nftables Documentation and HOWTO}, https://wiki.nftables.org/}

\begin{itemize}
\tightlist
\item
  Selezione e filtraggio dei pacchetti
\item
  NAT
\item
  Tracciamento delle connessioni
\item
  Accounting di rete.
\end{itemize}

Per quanto riguarda il filtraggio dei pacchetti, abbiamo a disposizione
5 hook differenti:

\begin{itemize}
\tightlist
\item
  \passthrough{\lstinline!NF\_INET\_PRE\_ROUTING!}: è chiamato nel
  metodo \passthrough{\lstinline!ip\_rcv()!} del kernel, quindi ancora
  prima di entrare nel sottosistema di routing.
\item
  \passthrough{\lstinline!NF\_INET\_LOCAL\_IN!}: è chiamato nel momento
  in cui il pacchetto è indirizzato al local host.
\item
  \passthrough{\lstinline!NF\_INET\_FORWARD!}: è chiamato dopo che è
  stato effettuato il lookup nel sottosistema di routing.
\item
  \passthrough{\lstinline!NF\_INET\_POST\_ROUTING!}: è chiamato per i
  pacchetti che devono essere inoltrati, pochi istanti prima che il
  pacchetto entri fisicamente nel cavo.
\item
  \passthrough{\lstinline!NF\_INET\_LOCAL\_OUT!}: è chiamato per i
  pacchetti creati nel local host, prima che escano.
\end{itemize}

Molte utility, comunemente usate, si basano sul sottosistema netfilter:

\begin{itemize}
\tightlist
\item
  \passthrough{\lstinline!iptables!}: agisce come front-end per
  netfilter ed implementa un \textbf{firewall}. Fornisce un'interfaccia
  per aggiungere, eliminare, regole netfilter, aggiungere tabelle,
  ecc\ldots{}
\item
  \passthrough{\lstinline!nftables!}: è sviluppato da \textbf{netfilter}
  (l'organizzazione dietro al sottosistema netfilter). È un rimpiazzo
  per iptables che è riconosciuto avere problemi di scalabilità e
  performance.
\end{itemize}

In \passthrough{\lstinline!nftables!}, utilizzato in BlueAgent, le
regole sono organizzate in tabelle e catene: ad esempio, è stata
definita una catena di tipo \passthrough{\lstinline!nat!} per gestire la
traduzione degli indirizzi di rete.

Spesso le connessioni a livello applicativo hanno una connessione
mantenuta nel tempo ed una serie di eventi che determinano la storia
della comunicazione. Il \textbf{Connection Tracking}
(\passthrough{\lstinline!conntrack!}) Layer permette al kernel di tenere
traccia delle sessioni create, è la struttura dati fondamentale per
implementare il NAT. \passthrough{\lstinline!conntrack!} mantiene per
ogni connessione che deve usufruire del NAT le traduzioni necessarie per
gli indirizzi IP.
